% ch10 -- Hyperbit vs. FFGFT: Structural Comparison
\label{ch:vergleich_matzke_ffgft}
% Source document: 326

\epigraph{Both calculations work out.
They only look elegant at different points.}{Doc.\,326}

\section{Starting Points in Comparison}

\begin{matzbox}[Matzke/Hyperbit]
\status{SETZUNG} $e_i^2 = +1$ (one axiom).\\
Borrowed bridges: Bekenstein, Compton$\,{=}\,$Schwarzschild, Landauer at $T_P$.\\
Result: Black hole physics from combinatorics, $\mbit = \mP\ln 2/(2\pi)$.
\end{matzbox}

\begin{ffgftbox}[FFGFT]
\status{SETZUNG} $\xi_{\text{par}} = 4/30000$, torus $T^4$, winding quantization.\\
No borrowed bridges --- but higher entry price.\\
Result: Bit energy is scale-sensitive, $\Ebit = \hbar c / L$. \status{K}
\end{ffgftbox}

\begin{laierspur}
Two hikers want to reach the same summit. One takes the short,
steep route --- he needs three ropes from others (Bekenstein, Compton, Landauer).
The other builds the entire path himself, pays the price upfront, and needs no help afterwards.
Both arrive. But only one has truly managed without foreign material.
\end{laierspur}

\section{Convergence: Where Both Agree}

Both frameworks agree: \status{B}
\begin{itemize}
  \item Geometry arises from distinguishability.
  \item Fermion generations depend on dimension-3 structures.
  \item Black holes mark an information threshold.
  \item The Planck regime is the natural limiting case.
\end{itemize}

\section{Divergence: The Landauer Point}

The essential difference: Matzke sets a universal bit mass;
FFGFT shows that $\Ebit$ depends on the scale. \status{K}

That is not only quantitative --- it is conceptual: If $\mbit$ is universal,
then all bits have the same physical status, regardless of whether they are encoded
in a DRAM, a spin or a photon.
If $\Ebit = \hbar c/L$, then \emph{there is no system-independent
bit mass} --- the number always depends on a physical carrier
with a physical scale.

\section{Cl(6) and GF(27)}

Doc.\,326 and 341 compare the algebraic structures:
$\mathrm{Cl}(6)$ with $\mathbb{Z}_3 \times \mathbb{Z}_3$ symmetry corresponds
in the FFGFT to the Galois structure $\mathrm{GF}(27) = \mathrm{GF}(3^3)$.
Three colors, three generations --- this three is geometrically grounded in both the hyperbit picture
and the FFGFT picture. \status{B}

Open Question \status{S}: Is $\mathrm{Cl}(6)$ a special case of the
FFGFT Galois structure on $T^4$, or are both independent paths to the
same physics?